\begin{textsample}{POS Dim 3 – rj – Score 26.00 – rj/rj\_inf\_cam\_exp\_14.txt}  \label{ex:f3_pos_004}
O Campo de Experiências Espaços, Tempos, Quantidades, Relações e Transformações Segundo a BNCC : As \textbf{crianças} vivem inseridas em espaços e tempos de diferentes dimensões, em um mundo constituído de fenômenos naturais e socioculturais.
Desde muito \textbf{pequenas}, elas procuram se situar em diversos espaços ( rua, bairro, cidade etc. ) e tempos ( \textbf{dia} e noite ; hoje, ontem e amanhã etc. ). \textbf{Demonstram} também \textbf{curiosidade} sobre o mundo físico ( seu próprio corpo, os fenômenos atmosféricos, os animais, as plantas, as transformações da natureza, os diferentes tipos de materiais e as possibilidades de sua manipulação etc. ) e o mundo sociocultural ( as relações de parentesco e sociais entre as pessoas que conhece ; como vivem e em que trabalham essas pessoas ; quais suas tradições e costumes ; a diversidade entre elas etc. ).
Além disso, nessas experiências e em muitas outras, as \textbf{crianças} também se \textbf{deparam}, frequentemente, com conhecimentos matemáticos ( contagem, ordenação, relações entre quantidades, dimensões, medidas, comparação de \textbf{pesos} e de comprimentos, avaliação de distâncias, reconhecimento de formas geométricas, conhecimento e reconhecimento de numerais cardinais e ordinais etc. ) que igualmente aguçam a \textbf{curiosidade}.
Portanto, a Educação \textbf{Infantil} precisa promover interações e \textbf{brincadeiras} nas quais as \textbf{crianças} possam fazer observações, \textbf{manipular} objetos, investigar e explorar seu entorno, levantar hipóteses e consultar fontes de informação para buscar respostas às suas \textbf{curiosidades} e indagações.
Assim, a \textbf{instituição} escolar está criando oportunidades para que as \textbf{crianças} ampliem seus conhecimentos do mundo físico e sociocultural e possam utilizá -los em seu cotidiano.
A \textbf{curiosidade}, o interesse e o prazer que as \textbf{crianças} apresentam nas situações em que podem criar cenários e enredos de histórias, fazer descobertas, resolver problemas do cotidiano, realizar uma tarefa com \textbf{colegas}, no campo de experiências ESPAÇOS, TEMPOS, QUANTIDADES, RELAÇÕES E TRANSFORMAÇÕES, nos levam a pensar em como lhes oferecer oportunidades para investigar as muitas questões que elas vão formulando acerca do mundo e de si mesmo e como nós, professores e professoras, podemos aprender mais sobre as \textbf{crianças} e suas formas de conhecer.
Temas como animais, plantas, sustentabilidade do ambiente, vida cotidiana, produção de bens e economia, nossa cidade, organizações sociais etc., e atividades que lidam com números, têm orientado o trabalho na Educação \textbf{Infantil}.
Estes e outros temas, no entanto, precisam ser tratados discutindo noções de espaço, de tempo, de quantidade, de relações e de transformações de elementos, quando se pretende motivar a \textbf{criança} a ter um olhar mais crítico e criativo do mundo, promovendo -lhe aprendizagens mais significativas.
Vivendo em uma aldeia, um sítio, uma fazenda ou um assentamento, as \textbf{crianças} desde \textbf{bebês} apreciam \textbf{brincar} com materiais da natureza.
Nas interações que estabelecem com seus familiares elas aprendem a reconhecer o cheiro da relva molhada, a chegada do momento de semear, de \textbf{colher}, o período de seca ou de chuva, os \textbf{sons} e as nuvens que anunciam a tempestade, os balidos dos carneiros ou da vaca, o comportamento das galinhas e das patas.
Imersas em um meio repleto de produtos da cultura, as \textbf{crianças} do campo e as \textbf{crianças} moradoras de zonas urbanas, ao \textbf{manipular} objetos e outros materiais, agem para entender seu funcionamento, para diferenciar suas características, cada vez mais com as \textbf{frequentes} perguntas “ como? ” e “ por quê? ” dirigidas a \textbf{parceiros} mais experientes.
Quanto tempo falta para o meu aniversário?
Por que quando minha avó era \textbf{criança} não havia televisão?
Por que alguns objetos afundam e outros não?
Por que existem alguns animais com penas e outros com pelos?
Quantas vezes um elefante é maior do que um cavalo?
Como estes doces podem ser distribuídos igualmente entre as \textbf{crianças}?
Que jogador de \textbf{futebol} fez mais gols na Copa?
Uma centopeia tem mais patas do que uma abelha?
Por outro lado, hoje todas as \textbf{crianças} observam situações em que os \textbf{adultos} lidam constantemente com pagamentos e trocos, calculam o tamanho de uma peça de \textbf{tecido} para poder fazer uma \textbf{vestimenta} ou o quanto de azulejo precisam comprar para finalizar uma casa, controlam o número de pessoas que estão presentes e o número de \textbf{dias} que faltam para uma determinada data etc., o que desperta nelas o \textbf{desejo} de se apropriar também desse saber.
Assim, elas gostam de perguntar a outras pessoas e de responder quantos anos têm, de \textbf{brincar} de telefonar fingindo discar ou digitar números, de trocar os canais da televisão, de recitar a seu modo a sucessão de números, de explorar as dimensões do espaço disponível no seu entorno etc.
A forma como os questionamentos das \textbf{crianças} da zona rural e da zona urbana são tratados pelos \textbf{adultos} próximos pode lhes ajudar ( ou não ) a observar regularidades e permanências, ou diversidades e mudanças na natureza e na vida social, a formular noções de espaço, de tempo e a fazer aproximações em torno da ideia de causalidade e transformação.
À medida que a \textbf{professora} considera as unidades de Educação \textbf{Infantil} como ambientes onde a \textbf{curiosidade} das \textbf{crianças} sobre o mundo físico e social pode alimentar a construção por elas de noções, comparações e implicações, ela as \textbf{ajuda} a construir explicações, conforme percebe seus \textbf{gestos}, sentimentos, intuições, seus motivos e sentidos pessoais nas respostas que elas dão.
Nesse processo procuram articular o modo como as \textbf{crianças} agem, sentem e pensam com os conhecimentos já disponíveis na cultura sobre cada objeto de conhecimento.

% matched lemmas: adulto, ajuda, bebê, brincadeira, brincar, colega, colher, criança, curiosidade, demonstrar, deparar, desejo, dia, frequente, futebol, gesto, infantil, instituição, manipular, parceiro, pequeno, peso, professora, som, tecido, vestimenta
\end{textsample}
